% !TEX root = ../main.tex
% =====================================================================
% 6. DISCUSSION AND LIMITATIONS -- page budget 0.70
%
% The trade-off paragraph the brief asks for, then the limitations,
% stated flatly. No new numbers are introduced here that Section 5 did
% not already report.
% =====================================================================

\section{Discussion and Limitations}
\label{sec:disc}

\paragraph{What the trade actually is.}
\Cref{sec:results} in one sentence: an order of magnitude of parameters and arithmetic, in exchange
for about two accuracy points on a broad archive, with the interpretation improved rather than
degraded where it can be scored. Whether that is a good trade is a property of the deployment, not of
the model. On a workstation it is a bad one. On a sensor node with a megabyte of flash it is the
difference between a method that can be considered and one that cannot, and the interpretation, which
is the reason to use MIL at all, arrives intact. Since the wide encoder recovers the \ucr{} deficit
(\Cref{sec:res:ucr}), what is being paid for is the size, and the payment can be adjusted by choosing
a width. The wider point: the tension between ``interpretable'' and ``small'' was never intrinsic to
MIL --- its head costs $0.25\,\%$ of the baseline --- but inherited from a backbone chosen when cost
was not a design variable.

\paragraph{Limitations.}
\textbf{No deployment}: we report parameter counts, analytically counted FLOPs and GPU timings; no
model was quantised, compiled or run on a microcontroller, and the int8 column of \Cref{tab:cost} is
arithmetic rather than measurement. \textbf{Activation memory did not shrink}: peak memory is worse
than the baseline's (\Cref{sec:res:cost}) because the encoder preserves length by design, so where
SRAM rather than flash binds this is the wrong design, and streaming or strided variants that keep a
per-step interpretation are the obvious next step. \textbf{Interpretation is scored on one dataset}:
\ndcg{} needs per-time-step ground truth, which only the synthetic \webtraffic{} generator supplies,
and that is also the screening set --- so the interpretation claim rests on one optimistic benchmark,
with AOPCR the weaker fallback (\Cref{sec:res:aopcr}). \textbf{Uneven seeds}: three models have three
seeds and the ablations are single-seed, read for shape against the spreads in \Cref{sec:setup}; a
full multi-seed sweep over 72 combinations was outside the available compute. \textbf{Training is
slower}: $117.3 \pm 14.0$\,s to converge on \webtraffic{} against $50.4 \pm 9.5$\,s for the baseline,
because depthwise convolutions are poorly served by GPU kernels tuned for dense ones. The saving is
at inference, where the target deployment spends its life, but the training cost is real.
